\documentclass[11pt,a4paper]{article}
\usepackage[margin=1in]{geometry}
\usepackage{hyperref}
\usepackage{enumitem}
\usepackage{graphicx}
\usepackage{xcolor}

% -----------------------------------------------------
% bvraghav-tiet-ucs503p-article.sty
% -----------------------------------------------------

% Variable: Lab Instructor
\makeatletter \newcommand{\BvrSetLabInstructor}[1]{%
  \gdef\Bvr@LabInstructor{#1}}
% default
\newcommand{\Bvr@LabInstructor}{Dr. Raghav %
  B. Venkataramaiyer}
% public accessor
\newcommand{\BvrLabInstructorValue}{\Bvr@LabInstructor}
\makeatother

% Add subtitle
\usepackage{titling}
% -- Set title bold
\pretitle{\begin{center}\bfseries\fontsize{18bp}{18bp}\selectfont}
\makeatletter
\newcommand{\BvrSubtitle}[1]{%
  \posttitle{%
    \par\end{center}
    \begin{center}\large#1\end{center}
    \vskip 0.5em}%
}
\makeatother

% Hints in the section title(s)
\newcommand{\BvrHint}[1]{%
  {\normalfont\normalsize\itshape\hfill #1}}

% -----------------------------------------------------

\title{AI Cost Autopilot: A Compiler-Aware Gateway for
  Token-Cost Optimisation in LLM Coding Agents}

\BvrSetLabInstructor{Dr. Raghav B. Venkataramaiyer}

\BvrSubtitle{%
  Project Proposal for UCS503P  \\ \bfseries
  Thapar Institute of Engineering and Technology%
}

\author{
  Agampreet Kaur, CSED%
  \thanks{Roll No: \texttt{1024240033}, %
    Email: \texttt{akaur11\_be24@thapar.edu}}
  \and
  Devansh Sharma, CSED%
  \thanks{Roll No: \texttt{1024240012}, %
    Email: \texttt{dsharma10\_be24@thapar.edu}}
  \and
  Furmaandeep Kaur, CSED%
  \thanks{Roll No: \texttt{1024240029}, %
    Email: \texttt{fsaini\_be24@thapar.edu}}
}

\begin{document}
\maketitle

\section{Higher-order goal}

This project aims to make AI-assisted software development economically
sustainable by reducing the computation paid for on every request to a Large
Language Model (LLM). While the topic is broader than any single development
team, the immediate engineering objective is to translate \emph{efficient
context usage} into a measurable, deployable piece of software: a gateway that
sits between a coding agent and the model it calls.

\section{Time-to-value}

\begin{itemize}[noitemsep]
  \item We shall deliver meaningful increments quickly: a transparent
        pass-through gateway first, then one optimisation at a time, each
        measurable on its own.
  \item We shall prioritise fast feedback over perfection, using weekly
        iterations with CI-backed automation from the first week to reduce
        release friction.
  \item Success is defined by what can be delivered and \emph{measured} in the
        first iteration, then improved --- every subsequent component must
        demonstrate a token reduction that the previous configuration did not
        already achieve.
\end{itemize}

\section{Problem Statement}

Developers increasingly work through LLM coding agents such as Cursor, Claude
Code and Codex. These tools operate by transmitting source code and
instructions to a hosted model, which bills per \emph{token}. As a codebase and
a team grow, both the size and the number of these requests grow with them.
Common pain points include:

\begin{itemize}[noitemsep]
  \item \textbf{Bulk context transmission:} agents attach large portions of a
        repository to every request, including files the task does not need.
  \item \textbf{Repetition:} semantically identical questions are re-asked and
        re-billed, because nothing remembers previous answers.
  \item \textbf{Budget unpredictability:} spend scales with usage in a way that
        teams cannot forecast or cap.
  \item \textbf{Exposure of private code:} more source leaves the developer's
        machine than the task requires.
  \item \textbf{Unsafe existing remedies:} generic prompt compressors treat code
        as plain text and may delete a line of logic the model needed, trading
        correctness for savings.
\end{itemize}

\textbf{Why this matters now (contextual relevance):} adoption of coding agents
is rising faster than any discipline for controlling their cost, and no standard
layer exists between the developer's tool and the model to manage it.

\section{Proposed Solution}

\subsection{Overview}
Build a \textbf{provider-agnostic gateway} that intercepts each request and
reduces its cost before forwarding it, with the following components:

\begin{itemize}[noitemsep]
  \item \textbf{Gateway:} an OpenAI/Anthropic-compatible endpoint, so an
        existing tool connects by changing a single base-URL setting.
  \item \textbf{Compiler-aware trimmer:} removes code the task does not need,
        using the structure of the code rather than its text.
  \item \textbf{Semantic cache:} reuses a stored answer when a new request means
        the same thing as an earlier one.
  \item \textbf{Autopilot policy:} selects, per request, the cheapest action that
        is still safe.
  \item \textbf{Metering:} records tokens before and after, cost, and cache
        hit-rate for every request.
\end{itemize}

\subsection{Core workflow}
\begin{enumerate}[noitemsep]
  \item The coding agent issues a request to the gateway instead of the provider.
  \item The gateway computes cheap signals (token count, presence of code) and
        the autopilot selects a plan.
  \item If a semantically equivalent request was answered before, the stored
        answer is returned without contacting the model.
  \item Otherwise the trimmer reduces the request, preserving the code the task
        depends on.
  \item The reduced request is forwarded to the provider; the answer is returned
        and recorded against the \emph{original} request.
\end{enumerate}

\subsection{Operational constraints for an educational setting}
\begin{itemize}[noitemsep]
  \item The system must run on student laptops: no GPU, no cluster.
  \item Every optimisation must degrade gracefully --- if an optional dependency
        is absent, the gateway continues to work as a pass-through.
  \item We avoid claims of novel research; the contribution is engineering
        judgement, measurement discipline and safety, not a new algorithm.
\end{itemize}

\section{Solution Approach}

\subsection{Context reduction}
Reduction is deterministic and structural, not statistical. Source files are
parsed into an Abstract Syntax Tree; a call graph is derived; the function the
task concerns and the functions it depends on are retained in full, while every
other definition is reduced to its signature. Because edits are made at syntactic
boundaries, the reduced file remains valid source. Where a dependency cannot be
resolved with confidence, the approximation deliberately retains \emph{more} than
necessary: an unnecessary retention costs savings, whereas an incorrect removal
costs correctness.

\subsection{System architecture}
A three-tier arrangement:
\begin{itemize}[noitemsep]
  \item \textbf{Interface:} an HTTP endpoint mirroring the provider API, so no
        client modification is required; streaming and tool-call structures pass
        through unchanged.
  \item \textbf{Optimisation layer:} the trimmer, the cache and the policy, each
        an independent module behind a fixed interface so they can be developed
        and tested in isolation.
  \item \textbf{Provider adapter:} forwards to the upstream model, with an
        offline mock so the whole pipeline can be exercised without an API key.
\end{itemize}

\subsection{CI/CD and fast delivery enabler}
Continuous integration runs the full test suite on every change, so a
modification to one module cannot silently break another. Every change reaches
the main branch through a reviewed pull request. Project documentation is built
and published automatically from the repository.

\section{Evaluation Criterion}

\subsection{Primary evaluation metric}
\textbf{Incremental token reduction at equal task outcome:} the percentage of
input tokens removed by our system \emph{beyond} what a naive alternative
already achieves, on requests where the agent still completes the task
correctly.

\begin{itemize}[noitemsep]
  \item Target: $\geq 30\%$ reduction against an untouched request, on
        code-heavy repetitive workloads.
  \item Attribution: measured with a tokeniser on the exact payloads sent, logged
        per request by the gateway itself.
\end{itemize}

\subsection{Secondary evaluation metrics}
\begin{itemize}[noitemsep]
  \item \textbf{Behavioural equivalence:} the agent reaches the same outcome with
        and without the gateway, over repeated runs. A saving that changes the
        result is counted as a failure, not a saving.
  \item \textbf{Cache hit-rate} and the proportion of cost avoided by reuse.
  \item \textbf{Added latency:} time introduced by the gateway per request.
  \item \textbf{Context exposure:} how much source code leaves the machine,
        relative to the unmodified agent.
\end{itemize}

\subsection{Validation plan}
A single fixed workload is run through progressively more capable
configurations, and the \emph{incremental} contribution of each is reported:
(i)~the untouched request, as a control; (ii)~trivial textual cleanup;
(iii)~generic text compression; (iv)~our compiler-aware reduction;
(v)~with semantic reuse enabled. This ablation is deliberate: a single
percentage quoted without a baseline cannot distinguish our contribution from
what a trivial transformation would have achieved anyway. Each row is reported
alongside its behavioural-equivalence result, so cost and correctness are never
presented separately.

\section{Scalability}

\begin{enumerate}[noitemsep]
  \item \textbf{Scalable (theoretically):} the gateway holds no per-request
        state, so instances may be added independently; cached entries are
        partitioned by request settings so lookups do not degrade as unrelated
        traffic grows; the store is addressed through an interface that admits a
        dedicated vector database.
  \item \textbf{Operationally deployable (practically):} it runs as a single
        local process for one developer, or as a shared instance for a team,
        where a common cache increases the proportion of reused answers.
\end{enumerate}

\section{Engine availability heuristic}

Mature components exist for every part of this system, so effort can be spent on
design and measurement rather than infrastructure:

\begin{itemize}[noitemsep]
  \item Web framework for an HTTP API, and an HTTP client for forwarding.
  \item Incremental parser with grammars for mainstream languages.
  \item Tokeniser for measurement, matching provider billing.
  \item Sentence-embedding model and a vector store for semantic reuse.
  \item Test framework, continuous-integration runner and documentation
        publishing.
\end{itemize}

\section{Project scope and deliverables}

\subsection{Initial deliverable}
\begin{itemize}[noitemsep]
  \item Gateway forwarding requests to a provider, with token and cost logging.
  \item Structural reduction for one language, preserving valid syntax.
  \item Exact-match reuse of previously answered requests.
  \item A policy selecting between reuse, reduction and pass-through.
  \item CI running the test suite on every change; published documentation.
\end{itemize}

\subsection{Subsequent deliverables}
\begin{itemize}[noitemsep]
  \item Dependency-aware retention across multiple files.
  \item Semantic reuse with a calibrated similarity threshold.
  \item A second language grammar.
  \item The ablation benchmark and the measured results it produces.
\end{itemize}

\section{Risks and mitigations}

\begin{itemize}[noitemsep]
  \item \textbf{Reduction removes needed context:} retain the task's
        dependencies in full, bias every approximation towards retention, and
        fall back to pass-through when a file cannot be parsed.
  \item \textbf{Reuse returns a plausible but wrong answer:} gate reuse behind a
        conservative similarity threshold calibrated on source code rather than
        prose, and never reuse across differing model or request settings.
  \item \textbf{Savings are indistinguishable from a trivial baseline:} report
        incremental gains against explicit baselines rather than a single
        headline figure.
  \item \textbf{Provider interfaces change:} isolate provider specifics behind a
        thin adapter.
  \item \textbf{Scope creep:} one language and one provider first; additional
        coverage only after the first iteration is measured.
\end{itemize}

\section{Summary}

This project proposes a deployable gateway that reduces the token cost of LLM
coding agents without changing developer workflow or agent behaviour. It meets
the course criteria by emphasising time-to-value through weekly measurable
increments, using CI/CD to support frequent safe releases, and defining
evaluation criteria that are both measurable and attributable --- specifically,
incremental token reduction reported against explicit baselines and paired with
behavioural equivalence. It remains focused on engineering workflow, safety and
measurement rather than research novelty.

\end{document}